\section*{CHAPTER 4. CONCLUSION}
\addcontentsline{toc}{section}{\numberline{}CHAPTER 4. CONCLUSION}

\setcounter{section}{4}
\setcounter{subsection}{0}
\setcounter{figure}{0}
\setcounter{table}{0}
\setcounter{equation}{0}

%============================================================%
\subsection{Summary of the Study}

This report presented a comparative study between Orthogonal Frequency
Division Multiplexing (OFDM) and Single-Carrier Frequency Division Multiple
Access (SC-FDMA), with emphasis on their Bit Error Rate (BER) performance
and Peak-to-Average Power Ratio (PAPR) characteristics.

Both transmission schemes were implemented in MATLAB under identical
simulation conditions, including:

\begin{itemize}
    \item 16-QAM modulation
    \item FFT size of 64
    \item 16 active subcarriers
    \item Cyclic Prefix length of 16
    \item Multipath Rayleigh fading channel
    \item Zero-Forcing frequency-domain equalization
\end{itemize}

A large-scale Monte Carlo simulation consisting of 200,000 transmitted
blocks was performed to ensure statistically reliable BER and PAPR
measurements.

%============================================================%
\subsection{Discussion of BER Performance}

Simulation results showed that both OFDM and SC-FDMA achieved reliable
communication performance over the multipath fading channel.

At low SNR values, OFDM demonstrated slightly better BER performance than
SC-FDMA. However, as the SNR increased, SC-FDMA gradually outperformed
OFDM and achieved lower BER at medium and high SNR regions.

This behavior can be explained by the DFT spreading operation used in
SC-FDMA. Since transmitted symbols are distributed across multiple
subcarriers, the system becomes more robust against severe fading on
individual subcarriers.

For example:

\begin{itemize}
    \item At 0 dB SNR:
    \[
    \mathrm{BER}_{\mathrm{OFDM}} = 1.859452 \times 10^{-1}
    \]
    
    \[
    \mathrm{BER}_{\mathrm{SC-FDMA}} = 2.249893 \times 10^{-1}
    \]

    \item At 20 dB SNR:
    \[
    \mathrm{BER}_{\mathrm{OFDM}} = 2.992188 \times 10^{-5}
    \]

    \[
    \mathrm{BER}_{\mathrm{SC-FDMA}} = 0
    \]
\end{itemize}

The BER curves also showed the expected monotonic improvement as SNR
increased, indicating that the system implementation and equalization
process operated correctly.

%============================================================%
\subsection{Discussion of PAPR Performance}

The PAPR analysis demonstrated one of the most important advantages of
SC-FDMA over OFDM.

The average measured PAPR values were:

\begin{equation}
\mathrm{PAPR}_{\mathrm{OFDM}} = 6.0139 \ \mathrm{dB}
\end{equation}

\begin{equation}
\mathrm{PAPR}_{\mathrm{SC-FDMA}} = 4.8222 \ \mathrm{dB}
\end{equation}

Therefore, SC-FDMA achieved approximately:

\begin{equation}
\Delta \mathrm{PAPR}
=
6.0139 - 4.8222
=
1.19 \ \mathrm{dB}
\end{equation}

of average PAPR reduction compared to OFDM.

The CCDF curves further confirmed that OFDM experiences significantly
higher probabilities of large instantaneous power peaks.

For example, at a PAPR threshold near 8 dB:

\begin{itemize}
    \item OFDM still exhibited noticeable probability of occurrence
    \item SC-FDMA probability had already approached zero
\end{itemize}

This lower PAPR property is highly important in practical wireless
transmitters because it improves power amplifier efficiency and reduces
nonlinear distortion.

%============================================================%
\subsection{Engineering Implications}

The results explain why OFDM and SC-FDMA are used differently in practical
cellular systems.

OFDM is highly suitable for downlink transmission because base stations
have:

\begin{itemize}
    \item Large power budgets
    \item More sophisticated hardware
    \item Less stringent amplifier efficiency constraints
\end{itemize}

In contrast, uplink mobile devices operate under limited battery capacity.
Therefore, reducing PAPR becomes critically important.

SC-FDMA preserves many OFDM advantages such as:

\begin{itemize}
    \item Frequency-domain equalization
    \item Robustness against multipath fading
    \item Flexible subcarrier allocation
\end{itemize}

while simultaneously improving transmitter power efficiency through its
single-carrier-like waveform characteristics.

This is the primary reason why SC-FDMA was selected for LTE uplink
transmission systems.

%============================================================%
\subsection{Limitations of the Study}

Although the implemented simulation successfully demonstrates the
fundamental differences between OFDM and SC-FDMA, several simplifying
assumptions were made:

\begin{itemize}
    \item Perfect synchronization was assumed
    \item Perfect channel estimation was assumed
    \item Only localized SC-FDMA mapping was considered
    \item A fixed modulation order (16-QAM) was used
    \item Nonlinear power amplifier effects were not modeled
\end{itemize}

In practical communication systems, these factors may further influence
system performance.

%============================================================%
\subsection{Future Work}

Several possible extensions can be explored in future studies:

\begin{itemize}
    \item Comparison between Localized and Interleaved SC-FDMA
    \item Investigation of higher-order modulation schemes
    \item Analysis under imperfect channel estimation
    \item Evaluation with nonlinear power amplifier models
    \item Comparison under realistic LTE or 5G numerologies
    \item Study of spectral efficiency and computational complexity
\end{itemize}

%============================================================%
\subsection{Final Remarks}

In conclusion, this study demonstrated the fundamental tradeoff between
OFDM and SC-FDMA.

OFDM provides excellent spectral efficiency and implementation simplicity,
but suffers from high PAPR due to the superposition of many independently
modulated subcarriers.

SC-FDMA mitigates this problem through DFT spreading, producing lower PAPR
and improved transmitter efficiency while retaining most OFDM advantages.

The simulation results verified that SC-FDMA achieves significantly lower
PAPR while maintaining competitive BER performance under multipath fading
conditions. These characteristics make SC-FDMA highly suitable for uplink
wireless communication systems where power efficiency is critical.